\begin{abstract}
A conformal prediction region promises to contain a vehicle's true future path a fixed fraction of the time. We ask
what happens to that promise when a trajectory forecaster is calibrated on one driving dataset and deployed, without
new labels, on another. Using AutoBot-Ego forecasters calibrated on Argoverse 2 and evaluated on nuScenes and on
Waymo Open Motion, and the reverse direction with nuScenes as the source, we measure a coverage gap of up to 23
percentage points at a 90\% nominal level. The gap is confirmed on models trained to match published state-of-the-art
accuracy and is pre-registered on the third dataset before it was touched. Stronger point prediction makes the gap
larger, not smaller. Reweighting the calibration set by an estimated covariate-shift likelihood ratio makes it worse,
and a label-free normalised score only partially repairs it. We give a training-conditional guarantee that holds
even when labels arrive scene by scene rather than one agent at a time, the realistic case for driving data, where a
naive per-agent guarantee is shown to fail outright. With 30 to 40 labelled target scenes it restores coverage close
to nominal at roughly 1.5 to 2 times the size of an oracle region, a result that replicates across both transfer
directions and all three datasets. We also localise the failure to turning manoeuvres and left-hand traffic, and
show that several measured input-level differences between the datasets explain only a minority of it. The results
argue that conformal guarantees for trajectory prediction should not be assumed to transfer across datasets without
an explicit, budgeted recalibration step.
\end{abstract}
